% Canonical LaTeX source; imported and revised from the legacy Markdown.
\chapter{Continuità}
\label{chap:8}

\begin{obiettivocapitolo}{verificare continuità, prolungamenti e parametri in funzioni definite a tratti}
\prioritaalta. Nei punti interni alle espressioni elementari la continuità segue dalle regole; il lavoro reale è nei punti di raccordo, negli estremi del dominio e nelle singolarità eliminabili.
\end{obiettivocapitolo}
\begin{proceduraesame}{continuità in un punto di raccordo}
\begin{enumerate}
\item Calcola $f(x_0)$ se definito.
\item Calcola separatamente $\lim_{x\to x_0^-}f(x)$ e $\lim_{x\to x_0^+}f(x)$ usando i rami corretti.
\item La continuità richiede l'uguaglianza dei due laterali e del valore $f(x_0)$.
\item Se compaiono parametri, risolvi il sistema ottenuto dalle uguaglianze necessarie.
\end{enumerate}
\end{proceduraesame}
\begin{proceduraesame}{prolungamento continuo}
Se $f$ non è definita in $x_0$ ma esiste finito $L=\lim_{x\to x_0}f(x)$, definisci $\widetilde f(x_0)=L$. Se il limite è infinito, diverso nei due lati o inesistente, il prolungamento continuo non è possibile.
\end{proceduraesame}
\begin{sceltametodo}{teoremi globali}
Usa Weierstrass per esistenza di massimo/minimo su compatto; il teorema degli zeri per garantire una radice da un cambio di segno; i valori intermedi per raggiungere un livello; Heine--Cantor per continuità uniforme su compatto.
\end{sceltametodo}
\begin{errorefrequente}{teorema degli zeri}
Il cambio di segno senza continuità non basta; la continuità senza cambio di segno non garantisce una radice. Il teorema garantisce esistenza, non il valore né l'unicità.
\end{errorefrequente}

Quando \(x_0\) è un punto di accumulazione del dominio, la continuità richiede tre elementi: \(f(x_0)\) deve essere definita, il limite per \(x\to x_0\) deve esistere ed essere finito, e deve coincidere con \(f(x_0)\). Nei punti isolati la continuità è automatica.

\section{Definizione e caratterizzazione mediante il limite}\label{definizione-e-caratterizzazione-mediante-il-limite}

Per \(f:D\to\mathbb R\) e \(x_0\in D\):

\[
f\text{ continua in }x_0
\]

\[
\Longleftrightarrow
\forall\varepsilon>0,\ \exists\delta>0,\ \forall x\in D:\
|x-x_0|<\delta
\Rightarrow
|f(x)-f(x_0)|<\varepsilon.
\]

Se \(x_0\in D\cap D'\):

\[
f\text{ continua in }x_0
\Longleftrightarrow
\lim_{\substack{x\to x_0\\x\in D}}f(x)=f(x_0).
\]

Caratterizzazione sequenziale:

\[
f\text{ continua in }x_0
\Longleftrightarrow
\forall(x_n)\subseteq D:\ x_n\to x_0
\Rightarrow
f(x_n)\to f(x_0).
\]

Agli estremi di un intervallo \([a,b]\):

\[
f\text{ continua in }a
\Longleftrightarrow
\lim_{x\to a^+}f(x)=f(a),
\qquad
f\text{ continua in }b
\Longleftrightarrow
\lim_{x\to b^-}f(x)=f(b).
\]

Se \(x_0\) è isolato in \(D\), \(f\) è continua in \(x_0\).

Approfondimento: \href{https://ingegnerismo.it/matematica/continuita/}{continuità}.

\section{Operazioni con funzioni continue}\label{operazioni-con-funzioni-continue}

Se \(f,g\) sono continue in \(x_0\):

\[
\alpha f+\beta g,
\qquad
fg,
\qquad
|f|,
\qquad
f^n\ (n\in\mathbb N_+)
\]

sono continue in \(x_0\) nei rispettivi domini.

Se \(g(x_0)\ne0\):

\[
\frac1g,
\qquad
\frac fg
\]

sono continue in \(x_0\).

\[
D_{f\pm g}=D_f\cap D_g,
\qquad
D_{fg}=D_f\cap D_g,
\]

\[
D_{f/g}=\{x\in D_f\cap D_g:g(x)\ne0\}.
\]

\[
\max(f,g)=\frac{f+g+|f-g|}{2},
\qquad
\min(f,g)=\frac{f+g-|f-g|}{2}.
\]

Se \(f\) è continua in \(x_0\) e \(g\) è continua in \(f(x_0)\):

\[
g\circ f\text{ è continua in }x_0,
\]

\[
D_{g\circ f}=\{x\in D_f:f(x)\in D_g\}.
\]

\[
\begin{array}{c|c}
\text{famiglia}&\text{dominio naturale di continuità}\\
\hline
\text{polinomi}&\mathbb R\\
\text{razionali}&\{Q\ne0\}\\
\sqrt[2m]{f}&\{f\ge0\}\\
\log_a f&\{f>0\},\ a>0,\ a\ne1\\
a^x&\mathbb R,\ a>0\\
\sin,\cos,\sinh,\cosh,\tanh&\mathbb R\\
\tan,\sec&\mathbb R\setminus\{\pi/2+k\pi\}\\
\cot,\csc&\mathbb R\setminus\{k\pi\}
\end{array}
\]

Approfondimento: \href{https://ingegnerismo.it/matematica/continuita/}{continuità}.

\section{Continuità dell'inversa monotona}\label{continuituxe0-dellinversa-monotona}

Se \(I\) è un intervallo e \(f:I\to\mathbb R\) è continua e strettamente monotona:

\[
J=f(I)\text{ è un intervallo},
\]

\[
f^{-1}:J\to I\text{ è continua e strettamente monotona}.
\]

\[
f\nearrow\Longrightarrow f^{-1}\nearrow,
\qquad
f\searrow\Longrightarrow f^{-1}\searrow.
\]

Approfondimento: \href{https://ingegnerismo.it/matematica/continuita-della-funzione-inversa/}{continuità della funzione inversa}, \href{https://ingegnerismo.it/matematica/funzione-inversa/}{funzione inversa}.

\section{Continuità uniforme e funzioni lipschitziane}\label{continuituxe0-uniforme-e-funzioni-lipschitziane}

La continuità uniforme è una proprietà dell'intero dominio, non di un singolo punto. Una funzione \(f:D\to\mathbb R\) è uniformemente continua su \(D\) se

\[
\forall\varepsilon>0\ \exists\delta>0\ \forall x,y\in D:
|x-y|<\delta
\Longrightarrow
|f(x)-f(y)|<\varepsilon.
\]

Negazione:

\[
\exists\varepsilon_0>0\ \forall\delta>0\ \exists x,y\in D:
|x-y|<\delta,
\qquad
|f(x)-f(y)|\ge\varepsilon_0.
\]

Condizione di Lipschitz:

\[
|f(x)-f(y)|\le L|x-y|
\quad\forall x,y\in D,
\qquad L\ge0,
\]

\[
f\text{ lipschitziana}
\Longrightarrow
f\text{ uniformemente continua}.
\]

\section{Tipi di discontinuità e prolungamento continuo}\label{tipi-di-discontinuituxe0-e-prolungamento-continuo}

Per \(x_0\in D\) punto di accumulazione del dominio da entrambi i lati:

\[
\begin{array}{c|c}
\text{tipo}&\text{condizione}\\
\hline
\text{eliminabile}&\lim_{x\to x_0}f(x)=L\in\mathbb R,\ f(x_0)\ne L\\
\text{a salto}&f(x_0^-),f(x_0^+)\in\mathbb R,\ f(x_0^-)\ne f(x_0^+)\\
\text{infinita}&f(x_0^-),f(x_0^+)\text{ esistono in }\overline{\mathbb R}
\text{ e almeno uno vale }\pm\infty\\
\text{oscillatoria o irregolare}&\text{almeno un limite laterale non esiste in }\overline{\mathbb R}
\end{array}
\]

Se \(x_0\notin D\), \(x_0\in D'\) e

\[
\lim_{\substack{x\to x_0\\x\in D}}f(x)=L\in\mathbb R,
\]

il prolungamento continuo è

\[
\widetilde f(x)=
\begin{cases}
f(x),&x\in D,\\
L,&x=x_0.
\end{cases}
\]

Approfondimento: \href{https://ingegnerismo.it/matematica/punti-di-discontinuita/}{punti di discontinuità}, \href{https://ingegnerismo.it/matematica/continuita/}{continuità}.

\section{Teorema di Weierstrass}\label{teorema-di-weierstrass}

\[
f\in C([a,b])
\Longrightarrow
\exists x_m,x_M\in[a,b],\ \forall x\in[a,b]:
f(x_m)\le f(x)\le f(x_M).
\]

\[
f(x_m)=\min_{[a,b]}f,
\qquad
f(x_M)=\max_{[a,b]}f.
\]

Approfondimento: \href{https://ingegnerismo.it/matematica/teorema-di-weierstrass/}{teorema di Weierstrass}.

\section{Teorema degli zeri}\label{teorema-degli-zeri}

\[
f\in C([a,b]),
\qquad
f(a)f(b)<0
\]

\[
\Longrightarrow
\exists c\in(a,b):\ f(c)=0.
\]

\[
f\in C([a,b]),
\qquad
f(a)f(b)<0,
\qquad
f\text{ strettamente monotona}
\Longrightarrow
\exists!c\in(a,b):\ f(c)=0.
\]

Approfondimento: \href{https://ingegnerismo.it/matematica/teorema-degli-zeri/}{teorema degli zeri}.

Esercizi svolti: \href{https://ingegnerismo.it/matematica/teorema-degli-zeri-esercizi/}{teorema degli zeri e bisezione}.

\section{Teorema dei valori intermedi}\label{teorema-dei-valori-intermedi}

\[
f\in C([a,b]),
\qquad
y\in[\min\{f(a),f(b)\},\max\{f(a),f(b)\}]
\]

\[
\Longrightarrow
\exists c\in[a,b]:\ f(c)=y.
\]

Equivalentemente, per ogni intervallo \(I\):

\[
f\in C(I)
\Longrightarrow
f(I)\text{ è un intervallo}.
\]

Approfondimento: \href{https://ingegnerismo.it/matematica/teorema-dei-valori-intermedi/}{teorema dei valori intermedi}.

\section{Teorema di Heine--Cantor}\label{teorema-di-heinecantor}

\[
f\in C([a,b])
\Longrightarrow
\forall\varepsilon>0,\ \exists\delta>0,\ \forall x,y\in[a,b]:
\]

\[
|x-y|<\delta
\Rightarrow
|f(x)-f(y)|<\varepsilon.
\]

Approfondimento: \href{https://ingegnerismo.it/matematica/teorema-di-heine-cantor/}{teorema di Heine--Cantor}.
